\chapter{Methodik}\label{ch:methodik}

In diesem Kapitel wird das Vorgehen zur Lösung der Problemstellung detailliert beschrieben.
Es ist das Herzstück der Arbeit, da es die eigene Leistung dokumentiert.
Die Beschreibung soll so konkret sein, dass Dritte die Arbeit nachvollziehen oder reproduzieren können.

Mögliche Inhalte:
\begin{itemize}
\item Formulierung der Problemstellung in mathematischer oder konzeptioneller Form.
\item Beschreibung des Lösungsansatzes (z. B. Algorithmus, Modell, Prototyp, Experimentdesign).
\item Datenbasis, Messverfahren oder Simulationsparameter.
\item Software- oder Systemarchitektur, Programmbibliotheken, verwendete Hardware.
\item Validierungskriterien und Messgrößen.
\end{itemize}

Verwenden Sie bei Bedarf Diagramme, Flussbilder oder Pseudocode, um komplexe Abläufe anschaulich darzustellen.



